\section{Defense: Topology-Aware DP Allocation}
\label{sec:defense}

The defense controls one knob per client: the noise scale $\sigma_i$
applied during DP-SGD. Standard practice sets $\sigma_i$ uniformly
across clients, implicitly assuming a symmetric topology. We allocate
$\sigma_i$ asymmetrically as a function of structural-leverage scores,
prove a per-client mutual-information bound that holds against the
adversary of \S\ref{sec:adv-model}, and show the allocation is balanced
min-max optimal: it strictly dominates uniform DP-SGD whenever the
leverage scores are non-uniform across clients.

We label the allocation \emph{Fulcrum}: per-client privacy bounds are
equilibrated across the asymmetric leverage profile at a fixed utility
budget, as a lever balances unequal loads about its fulcrum.

\subsection{Notation and Setup}
\label{sec:defense-setup}

The world prior $\mathbb{P}_{\mathrm{world}}$ generates client datasets
$D = (D_1,\ldots,D_n)$ with class distributions $\Delta_i = \delta(D_i)$
and $p_i = \Delta_i(\Cs)$. DP-SGD~\cite{abadi2016dpsgd} produces, per
round and per client,
\begin{equation}
\theta_i^{(t)} = \theta^{(t-1)} - \eta \cdot \tilde g_i^{(t)},
\quad
\tilde g_i^{(t)} = \frac{1}{|B|}\sum_{x \in B_i^{(t)}} \mathrm{clip}\big(\nabla\ell(\theta^{(t-1)},x), C\big) + \xi_i^{(t)},
\label{eq:dpsgd}
\end{equation}
with $\xi_i^{(t)} \sim \mathcal{N}(0, \sigma_i^2 C^2 I)$. The adversary
observes $\Theta = \{\theta_i^{(t)}\}_{i,t}$ and outputs estimate
$\hat p_i$ via deterministic $f$ as in \S\ref{sec:attack}. We let
$\Tmax$ denote the observation-window length: the number of rounds
visible to the adversary, capped if a windowing/re-keying defense is in
place.

\paragraph{Prior family.} We constrain analysis to a family of priors
$\FAM_{\TOPO,\omega}$ over $(\Delta_1,\ldots,\Delta_n)$ that
(a)~factorise over connected components of $\TOPO$, (b)~within each
component satisfy a Markov property with respect to $\TOPO$, and
(c)~have bounded second moments. These conditions are standard in
graphical-model analyses and are satisfied by the realistic deployment
priors of \S\ref{sec:experiments}; condition (c) ensures that supremum
quantities below are finite.

\subsection{Structural Leverage}
\label{sec:defense-leverage}

\begin{definition}[Structural leverage]
\label{def:leverage}
For client $i$ in deployment $(\TOPO,\omega)$, the structural leverage is
\[
\ell_i^\circ \;:=\; \sup_{\mathbb{P} \in \FAM_{\TOPO,\omega}}
                     I_{\mathbb{P}}\big(p_i;\, D_{-i}\big).
\]
\end{definition}

The supremum captures the worst-case prior coupling between $p_i$ and
the rest of the federation's data $D_{-i}$, given the deployment
structure. It is a deterministic function of $(\TOPO,\omega,i)$, finite
under condition (c), and computable in closed form for the practical
proxies of \S\ref{sec:defense-proxies}.

\subsection{Per-Client Conditional MI Bound}
\label{sec:defense-thm1}

\begin{theorem}[Per-client conditional MI bound]
\label{thm:per-client}
Under (IA1)--(IA2), for any deterministic adversary strategy $f$,
any prior $\mathbb{P} \in \FAM_{\TOPO,\omega}$, and any client $i$:
\begin{equation}
I_{\mathbb{P}}\big(p_i;\,\hat p_i \;\big|\; \TOPO,\omega,\{\sigma_j\}\big)
\;\leq\;
\underbrace{\frac{\Tmax C^2}{2\sigma_i^2 |B|^2}}_{\text{controllable: DP-SGD on client } i}
\;+\;
\underbrace{\ell_i^\circ}_{\text{uncontrollable: prior coupling}}.
\label{eq:thm1}
\end{equation}
\end{theorem}

\begin{proof}[Proof sketch]
By data processing $\hat p_i = f(\Theta) \Rightarrow I(p_i;\hat p_i) \leq I(p_i;\Theta)$.
The conditioning identity with $D_{-i}$ and dropping the non-negative
subtractive term yields
$I(p_i;\Theta) \leq I(p_i;\Theta \mid D_{-i}) + I(p_i;D_{-i})$. The first
term reduces to $I(D_i;\Theta_i \mid D_{-i})$ under (IA2) and is bounded
by $\Tmax C^2/(2\sigma_i^2 |B|^2)$ via Mironov's RDP composition
\cite{mironov2017rdp} applied to the Gaussian
mechanism~\cite{abadi2016dpsgd}, converted to mutual information via
Cuff and Yu's max-KL bound~\cite{cuff2016mi}. The second term is bounded
by $\ell_i^\circ$ by definition. Full proof in Appendix~\ref{app:thm1}.
\end{proof}

The additive structure is the key feature of the bound. The first term is
\emph{controllable}: it shrinks as the per-client noise $\sigma_i$
grows. The second term is \emph{uncontrollable}: it is a property of the
prior and is not reduced by DP-SGD noise on client $i$. The allocation
of \S\ref{sec:defense-thm2} manages the controllable term given that
the uncontrollable floor is fixed.

\subsection{Balanced Min-Max Optimal Allocation}
\label{sec:defense-thm2}

Let $a := \Tmax C^2 / (2|B|^2)$. The optimisation problem is
\begin{equation}
\min_{\{\sigma_i^2 > 0\}}\; \max_i \big[\, a / \sigma_i^2 + \ell_i^\circ\, \big]
\quad \text{subject to} \quad
\sum_{i=1}^n \sigma_i^2 \leq U,
\label{eq:opt}
\end{equation}
where the budget $U$ aggregates per-client noise variance.
Proposition~\ref{prop:utility} below justifies treating $U$ as a
proxy for utility cost.

\begin{theorem}[Balanced min-max allocation]
\label{thm:allocation}
For any utility budget $U > 0$, the optimal allocation solving
\eqref{eq:opt} is
\begin{equation}
\sigma_i^{*\,2} \;=\; \frac{a}{K^\star - \ell_i^\circ},
\label{eq:sigma-star}
\end{equation}
where $K^\star$ is the unique solution to
\begin{equation}
\sum_{i=1}^n \frac{a}{K^\star - \ell_i^\circ} = U,
\quad
K^\star > \max_i \ell_i^\circ.
\label{eq:budget-eq}
\end{equation}
The achieved worst-case bound is $\max_i I = K^\star$, balanced uniformly
across clients.
\end{theorem}

\begin{proof}[Proof sketch]
Reformulate \eqref{eq:opt} via slack $K$:
$\min K$ subject to $a/\sigma_i^2 + \ell_i^\circ \leq K$ for all $i$
and $\sum_i \sigma_i^2 \leq U$. The objective is linear, the constraints
are convex (since $a/x$ is convex on $x>0$), and Slater's condition
holds at $\sigma_i^2 = U/(2n)$. KKT therefore gives
$\sigma_i^{*\,2} = a/(K^\star - \ell_i^\circ)$ at the active-budget
optimum; uniqueness of $K^\star$ follows from strict monotonicity of
the budget equation in $K$. Full proof in Appendix~\ref{app:thm2}.
\end{proof}

\begin{corollary}[Strict improvement over uniform allocation]
\label{cor:strict}
Under uniform allocation $\sigma_i^2 = U/n$, the worst-case bound is
$K_{\mathrm{uniform}} = an/U + \max_i \ell_i^\circ$. The balanced allocation of
Theorem~\ref{thm:allocation} satisfies
\(K^\star \leq K_{\mathrm{uniform}}\), with equality if and only if all
$\ell_i^\circ$ are equal.
\end{corollary}

By Corollary~\ref{cor:strict}, the topology-aware allocation helps
exactly when it should. Symmetric topologies (ring, complete, balanced hierarchical) produce
uniform leverage and the allocation reduces to uniform DP-SGD; this is
the appropriate behaviour, not a flaw, since there is no asymmetry to
exploit. Asymmetric topologies (line, hierarchical with uneven groups,
star) and deployments with heterogeneous client dataset sizes both
produce non-uniform leverage and admit a strict gap.

\subsection{Practical Leverage Proxies}
\label{sec:defense-proxies}

The abstract leverage $\ell_i^\circ$ is well-defined but not always
analytically computable. We instantiate three principled proxies, each
appropriate to a different deployment regime.

\begin{corollary}[SBM proxy, hierarchical FL]
\label{cor:sbm}
Under a stochastic block model prior on $(\Delta_1,\ldots,\Delta_n)$
with $\Delta_i \sim_{\text{iid}} \mathbb{P}_{\Phi_{\omega_i}}$ given
finite-entropy block parameter $\Phi$,
\(\ell_i^\circ \leq H(\Phi_{\omega_i})\),
asymptotically independent of block size; for finite blocks,
$\ell_i^\circ$ is approximated by the group size $|G_i|$ scaled by a
prior-dependent constant. This motivates the proxy
$\ell_i^{\mathrm{org}} \propto |G_i|$ for hierarchical FL with uneven
organisational groups (\eg Briggs et al.\ clustering~\cite{briggs2020hierarchical}).
\end{corollary}

\begin{corollary}[Degree heuristic, decentralised FL]
\label{cor:degree}
Under bounded-degree Markov priors with a regularity assumption on
neighborhood-conditional entropy, leverage admits an upper bound
proportional to graph degree. We adopt the proxy
$\ell_i^{\mathrm{deg}} \propto \deg_{\TOPO}(i)$ as a heuristic. The
heuristic is empirically validated in \S\ref{sec:experiments}; we do not
claim a tight proof. Appropriate for decentralised peer-to-peer
deployments~\cite{roy2019braintorrent,hegedus2021decentralized,beilharz2021dag}
with non-uniform topology degree.
\end{corollary}

\begin{corollary}[FedAvg-influence proxy, cross-silo FL]
\label{cor:influence}
In FedAvg with size-weighted aggregation
$\theta^{(t+1)} = \sum_i \frac{|\Dset_i|}{\sum_j |\Dset_j|} \theta_i^{(t+1)}$,
each client's update is weighted proportionally to its dataset size.
We adopt the proxy $\ell_i^{\mathrm{ds}} \propto |\Dset_i|/\bar{|\Dset|}$
(mean-normalised) as a heuristic for cross-silo settings where topology
is symmetric (\eg ring, complete) but client dataset sizes vary
substantially, as in the typical FLamby regime~\cite{flamby2022}. Like
Corollary~\ref{cor:degree}, this is empirically validated rather than
proven.
\end{corollary}

The three proxies are not mutually exclusive: a deployment can blend
them by passing per-client weight vectors that combine size, degree,
and group size, normalised to mean 1. We treat the choice of proxy as a
deployment design parameter and report sensitivity to it as part of
\S\ref{sec:experiments-ablations}.

\subsection{Utility Cost}
\label{sec:defense-utility}

\begin{proposition}[Utility cost]
\label{prop:utility}
Under standard FL convergence assumptions ($L$-smooth loss, bounded
local variance, bounded heterogeneity, full participation), the
optimisation error after $\Tmax$ rounds with per-client noise variances
$\{\sigma_i^2\}$ admits the decomposition
$f_T(\Tmax) + f_\sigma\big(\sum_i \sigma_i^2\big)$,
where $f_T$ is decreasing in $\Tmax$ and $f_\sigma$ is increasing in
its argument. The decomposition specialises to
$O(\log\Tmax / (\mu \Tmax)) + O(d\,V/\mu)$ for $\mu$-strongly
convex losses (cf.\ Wei et~al.~\cite{wei2020nbafl} adapted to the
per-client noise setting) and to $O(1/\sqrt{\Tmax}) + O(\sqrt{dV})$ for
non-convex smooth losses; here $V = (C^2 / |B|^2 n^2) \sum_i \sigma_i^2$
is the aggregated noise variance per coordinate and $d$ is the
parameter dimension. Convergence machinery follows
McMahan et~al.~\cite{mcmahan2018dpfedavg} for per-client noise and
Li et~al.~\cite{li2020fedavg} for the underlying non-IID FedAvg rate.
\end{proposition}

The proposition justifies treating $U = \sum_i \sigma_i^2$ as a
utility budget proxy and motivates the constraint set in
\eqref{eq:opt}. Concretely, increasing $U$ trades privacy (looser
$K^\star$) for utility (smaller convergence error). The Pareto frontier
of \S\ref{sec:experiments} sweeps $U$ to map this trade-off.

\subsection{When the Defense Matters}
\label{sec:defense-regimes}

The defense provides marginal benefit \emph{exactly when} the topology
or organisational structure produces non-uniform leverage. The relevant
regimes are summarised in Table~\ref{tab:regimes}.

\begin{table}[h]
\centering
\caption{Conditions under which the topology-aware allocation strictly
improves over uniform DP-SGD, with measured $K^\star$ gaps.
Gaps are $K_{\mathrm{uniform}} - K^\star$ at $\eta=1$ on Setting~C
($n=50$, $\Tmax=100$, $|B|=64$, $C=1.0$, $U=0.5$); the star gap
saturates near the theoretical asymptote $an/U = 1.2207$.}
\label{tab:regimes}
\begin{tabular}{@{}lllc@{}}
\toprule
Topology / structure & Leverage uniform? & Strict gap over uniform? & Measured gap (nats) \\
\midrule
Ring (uniform degree) & Yes & No (degenerate; \cf~Cor.~\ref{cor:strict}) & $0.0000$ \\
Complete (uniform degree) & Yes & No (degenerate) & --- \\
Line (asymmetric endpoints) & No & Small & $0.0242$ \\
Hierarchical with balanced groups & Yes & No (degenerate) & --- \\
Hierarchical with uneven groups\footnotemark[1] & No & Moderate & $0.4555$ \\
Star (1 hub + $n-1$ leaves) & No & Large; saturates at $\sim an/U$ & $1.1958$ \\
Cross-silo with heterogeneous $|\Dset_i|$\footnotemark[2] & No & Strong (\cf~Cor.~\ref{cor:influence}) & $0.6280$ \\
\bottomrule
\end{tabular}
\footnotetext[1]{Region sizes $[20,15,10,3,2]$; see \S\ref{sec:experiments-eta}.}
\footnotetext[2]{Setting~B (Fed-Heart-Disease, native FLamby sites
$[199, 172, 30, 85]$); maximum gap across the
$U \times \Tmax$ Pareto sweep is reported.
See \S\ref{sec:experiments-pareto}.}
\end{table}
